% Esqueleto didático adaptado do pré-projeto fornecido pelo aluno.
\section{Introdução}

A realização de viagens institucionais constitui uma atividade necessária para o funcionamento de diversas organizações públicas e entidades que executam atividades de interesse público. Deslocamentos aéreos são frequentemente utilizados para viabilizar reuniões técnicas, acompanhamento de projetos, ações de monitoramento, atividades de avaliação, participação em eventos, capacitações, articulação com parceiros e execução de agendas institucionais em diferentes unidades da federação. Nesse contexto, a aquisição de passagens aéreas não representa apenas uma rotina administrativa, mas uma decisão de gestão com impacto direto sobre a eficiência do gasto, a previsibilidade orçamentária e a capacidade de execução das atividades finalísticas.

O problema central não está apenas em saber se comprar com maior antecedência tende, em média, a estar associado a menores valores. A questão metodológica mais relevante é que, quando se analisam apenas passagens efetivamente compradas, observa-se somente o preço realizado em uma determinada data de aquisição. Não se observa, para aquela mesma passagem, qual teria sido o preço caso a compra tivesse ocorrido em outro momento. Assim, uma passagem comprada com 15 dias de antecedência não pode ser comparada diretamente com outra passagem comprada com 21 dias de antecedência como se ambas representassem o mesmo contrafactual, ainda que tenham origem e destino semelhantes. Para sustentar conclusões mais robustas sobre a variação de preços ao longo da antecedência, seria necessário observar preços sucessivos para a mesma rota, mesma data de viagem e condições de busca controladas, ou utilizar bases complementares que ampliem a validade da análise.
